\section{Síntesis y ampliación}

\subsection{Repaso de las ideas centrales}

En esta entrevista, Demis Hassabis presenta un panorama completo en torno a la definición de la AGI, su cronograma, su impacto y sus riesgos. A continuación, un resumen condensado de las ideas centrales:

\begin{table}[H]
\centering
\begin{tabular}{p{0.25\textwidth} p{0.65\textwidth}}
\toprule
\textbf{Tema} & \textbf{Juicio central} \\
\midrule
Cronograma de la AGI & Es muy probable que se logre dentro de cinco años \\
Magnitud del impacto de la AGI & Diez veces la Revolución Industrial, a diez veces la velocidad \\
Estado actual de la IA & «Inteligencia desigual»; la fragmentación es la causa raíz \\
Cuatro capacidades ausentes & Aprendizaje continuo, arquitectura de memoria, planeación jerárquica, razonamiento consistente \\
Leyes de escalamiento & Rendimientos decrecientes, pero no agotadas; la competencia se desplaza hacia la innovación algorítmica \\
Perspectiva de los LLM & No se volverán una mercancía; son un componente básico de la AGI \\
Brecha entre las compañías de vanguardia & Se está ampliando, no reduciendo \\
Estrategia de DeepMind & Integración de tres en uno, con tres capas de modelos: de vanguardia, de código abierto y especializados por dominio \\
Descubrimiento de fármacos & Resultados transformadores en 5 a 10 años, con la meta de vencer al cáncer \\
Riesgo de seguridad prioritario & El uso malicioso, por encima del alineamiento técnico \\
Política preferida & Establecer un equivalente de la OIEA para la IA \\
Energía & La IA es tanto causante como solucionadora del problema \\
Sesgo de percepción & Sobrestimada en el corto plazo, gravemente subestimada en el largo plazo \\
\bottomrule
\end{tabular}
\caption{Panorama de las ideas centrales de Hassabis}
\end{table}

\subsection{Citas clave}

Las citas directas más contundentes de esta entrevista:

\begin{enumerate}
    \item \textit{«A veces cuantifico la llegada de la AGI como diez veces la Revolución Industrial, ocurriendo a diez veces la velocidad.»} —sobre la magnitud del impacto de la AGI
    \item \textit{«Preguntar de una manera produce resultados asombrosos; cambiar la formulación hace que falle en lo más elemental.»} —sobre la «inteligencia desigual»
    \item \textit{«El 90\% de los avances que sostienen a la IA moderna fueron logrados por Google Brain, Research o DeepMind.»} —sobre la contribución histórica de DeepMind
    \item \textit{«Es muy probable que la AGI se logre dentro de los próximos cinco años.»} —sobre el cronograma, en continuidad con la predicción de 20 años que hizo Shane Legg en 2010
    \item \textit{«La IA está sobrevalorada en el corto plazo y gravemente subestimada en el largo plazo.»} —sobre el sesgo de percepción
\end{enumerate}

\subsection{Conclusiones más prácticas para el lector}

\begin{enumerate}
    \item \textbf{No subestime el cronograma de la AGI}: cinco años no es una exageración; Hassabis viene apostando por esta dirección desde 2010 y, hasta ahora, se ha mantenido en el rumbo correcto
    \item \textbf{Preste atención a la capacidad de adaptación institucional}: cuanto más rápidos son los avances tecnológicos, más se acentúan los cuellos de botella institucionales. La reforma regulatoria, la transformación educativa y el rediseño de los sistemas de seguridad social deben iniciarse con antelación
    \item \textbf{La seguridad de la IA requiere colaboración internacional}: no es un problema que pueda resolver una sola compañía o un solo país; requiere un mecanismo global similar a la OIEA
    \item \textbf{Pragmatismo en el corto plazo, respeto en el largo plazo}: no se deje arrastrar por el exceso de entusiasmo, pero tampoco ignore la profunda transformación de largo plazo por una «desilusión con la IA» de corto plazo
\end{enumerate}

\subsection{Lecturas adicionales}

\begin{itemize}
    \item Demis Hassabis et al., \textit{AlphaFold: Protein Structure Prediction with Deep Learning}, Nature, 2021
    \item Demis Hassabis et al., \textit{Mastering the Game of Go with Deep Neural Networks and Tree Search}, Nature, 2016
    \item Shane Legg, \textit{Machine Super Intelligence}, PhD Thesis, 2008
    \item Google DeepMind, \textit{Gemini Technical Report}, 2024
    \item Rich Sutton, \textit{The Bitter Lesson}, 2019
    \item Dario Amodei, \textit{Machines of Loving Grace}, 2024
\end{itemize}

\end{document}
